% cover_letter.tex -- the cover letter for the submission of "Higgs as a Composite, and Quantum Gravity" to Physical Review D.
% PROJECT spacetimes, folder Revision_SM_Higgs/paper_prd.  Generated from src/cover_letter.template.tex by
% scripts/build_paper.py, which inserts the numbers (the same placeholders as in the article, from the data files).
\documentclass[11pt]{article}
\usepackage[margin=2.5cm]{geometry}
\usepackage{amsmath,amssymb}
\usepackage[colorlinks=true,linkcolor=blue,urlcolor=blue]{hyperref}
\hypersetup{pdftitle={Cover letter: Higgs as a Composite, and Quantum Gravity},pdfauthor={Patrick Lee Nash}}
\setlength{\parindent}{0pt}
\setlength{\parskip}{6pt}
\begin{document}

\begin{flushright}
Patrick Lee Nash, Ph.D.\\
PROJECT spacetimes\\
November 21, 2026
\end{flushright}

The Editors\\
Physical Review D

\textbf{Submission of a Regular Article: ``Higgs as a Composite, and Quantum Gravity''}

Dear Editors,

I submit the manuscript ``Higgs as a Composite, and Quantum Gravity'' for consideration as a Regular Article in
Physical Review D, together with its Supplemental Material.

\textbf{Context.} In an earlier manuscript (``The Lorentz Super-Loop of the Bi-Octonions: Non-Associative Counterexamples
to the No-Go Theorems and a Spin-Zero Spinor Superfluid'', dated November 16, 2026, unpublished) I showed that the unit
bi-octonions form a Moufang loop that contains $SL(2,\mathbb C)$ together with eight spinor charges, a Weyl spinor that is
an isospin doublet, and that a spin-zero superfluid of these charges forms isospin-triplet pairs. The present manuscript
asks whether such a superfluid can be the Higgs sector of a gauge theory. The spinor field is a DiracLike spinor (a Dirac
spinor of $SL(2,\mathbb C)$ with electric charge zero and no electroweak quantum numbers), and its gauge group
$SL(2,\mathbb C)\times SU(2)_{\rm v}\times U(1)^{\rm F}$ is a new group, not the electroweak group. The manuscript solves the
superfluid in five mean-field methods (Hubbard-Stratonovich with three decouplings, Kohn-Sham, and
Cornwall-Jackiw-Tomboulis), derives the effective Lagrangians of the composite order parameter, and asks whether the
superfluid is the Higgs field of its own gauge group and whether it could be that of the Standard Model.

\textbf{Main findings.}
\begin{enumerate}
\item In every method the condensate is the Higgs field of the new group: the isospin bosons $W^\pm$ and the boson $Z$ of
the fermion number become massive, and exactly one gauge boson stays massless. The massless gauge bosons correspond one to
one to the generators that leave the condensate unchanged. The massless boson, which I name the zeugon, is a new gauge
boson, not the photon: it couples to the unbroken isospin charge, it is the only gauge boson that attracts the two members
of a pair, and, because the surviving group is $SL(2,\mathbb C)\times O(2)$, its charge conjugation is itself a gauge
symmetry, the structure of Alice electrodynamics.
\item The amplitude mode is a composite scalar bound below the threshold of two quasi-particles; the Higgs field of the
model is composite by construction (the compositeness condition of Bardeen, Hill and Lindner).
\item From the static pair response of a moving pair I compute the gradient coefficient of the order parameter at the
critical temperature and find Ginzburg numbers ${\rm Gi}={{GI_RANGE}}$: the critical region is wider than the critical
temperature itself, and the mean-field transition, on which all five methods agree, is not reliable. The calculation also
shows that a cutoff on the momenta of the two constituents destroys the gradient expansion (the coefficient of
$|\vec\Phi|^2$ acquires a term linear in the momentum of the pair), while a cutoff on their relative momentum does not.
\item The amplitude mode is a broad resonance: once it can decay into the longitudinal gauge bosons its tree-level width is
{{WIDTH_RANGE}} times its mass, and its mass is ${{MHV_RANGE}}$ times the scale of the breaking. Because the DiracLike spinor
carries no electroweak quantum numbers, the condensate is a singlet of the electroweak group and cannot be the Higgs field
of the Standard Model, whatever its numbers.
\end{enumerate}

\textbf{Why Physical Review D.} The manuscript concerns a Higgs sector, composite Higgs bosons and dynamical symmetry
breaking, a new massless gauge boson with a local charge conjugation, and the finite-temperature transition of a gauge
theory. Its treatment of the gradient coefficient and of the regularization of a moving pair applies to any
model in which a four-fermion interaction condenses a fermion field.

\textbf{On the title and the gravitational sector.} The title names quantum gravity because the gauge group of the model
contains $SL(2,\mathbb C)$ as a gauge factor, and Sec.~VIII of the manuscript outlines how E.~A.~Lord's gauge theory of the
local change of spinor representation will add the gravitational sector. The abstract and that section state that this
sector is a program and that no part of it is computed in this manuscript.

\textbf{Reproducibility.} Every number of the manuscript is produced by two independent programs (one in Rust, one in
Python with numpy and LAPACK); the statements of Secs.~III to VI, except the mass matrix of the gauge bosons and the
time-dependent calculations, are checked a third time by a notebook in the Wolfram Language. The Supplemental Material
gives the derivations, the complete tables, every check with its deviation and tolerance, and the SHA-256 checksums of the
data files and programs, which are openly available in a public repository (the Data Availability statement of the
manuscript).

\textbf{Submission history.} This manuscript has not been published and is not under consideration for publication
elsewhere; it has not been submitted to a journal before.

\textbf{Suggested referees.} Suggestions of referees, by area of expertise, are given in the editorial information of
the submission.

Thank you for your consideration.

Sincerely,

Patrick Lee Nash, Ph.D.

\end{document}
